\begin{ntlemma}[Author Lemma16]
Let $\ell$ be the largest integer such that $\operatorname{cover}([n],\mathcal S)\geq4\ell$. Then, $\Delta>\ell$ with high probability.
\end{ntlemma}
\begin{proof}
Our strategy will be as follows. We say that the choice of $T$ is bad if, for that choice of $T$, $\Delta\leq\ell$. We will then upper bound the number of bad $T$ by showing such $T$ have short descriptions.
Fix some bad $T$. Then $\ell\geq\Delta=\CCirc(T\bullet g)-k$, so $\CCirc(T\bullet g)\leq\ell+k$. Since there is a circuit $C$ computing $T\bullet g$ using at most $\ell+k$ gates and $k$ of the gates in $C$ must compute the unique components of $g$ (using Proposition14), it follows that there is a circuit $D$ that takes $(\log(m)+k)$-inputs and has at most $\ell$ gates such that
\[D(x,g_1(x),\ldots,g_k(x))=T(x)\]
for all $x\in\{0,1\}^{\log m}$ where $g_1,\ldots,g_k$ are the unique components of $g$. Moreover, since $D$ has only $\ell$ gates of fan-in 2, it uses at most $2\ell$ of the components of $g$ in the circuit. Thus, after a possible relabeling of $g_1,\ldots,g_k$, we can assume $D$ takes at most $(\log(m)+2\ell)$-inputs and that
\[D(x,g_1(x),\ldots,g_{2\ell}(x))=T(x).\]
Hence, to describe $T$, we just need to have a description for $D$ as well as a description for $g_1,\ldots,g_{2\ell}$. Indeed, this will be the last step in our eventual description of $T$. We present the eventual description now so as to guide the reader. Our proof will subsequently proceed working through this description from bottom to top.
\begin{enumerate}
\item Given
\begin{itemize}
\item a subset $J\subseteq[n]$ of size at most $\leq n(1-\frac1{2r})$
\item for each $j\in J$ a partial truth table encoding $(j,V_j)\in[n]\times\{0,1\}^{m/n+1}$
\item $r$-bounded subsets $S_1,\ldots,S_{2\ell}\subseteq[n]$ whose union is $J$,
\item gate numbers $u_1,\ldots,u_{2\ell}\in[2m\log m]$,
\item and a circuit $D$ of size $\ell$ with $(n+2\ell)$ inputs
\end{itemize}
\item For $j\in J$, let $T\langle P_j^{m,n}\rangle$ be the function whose values on $P_j^{m,n}$ in lexicographic order are given by the binary string $V_j$ and is zero everywhere else
\item For $i\in[2\ell]$, let $T\langle S_i^m\rangle=\bigvee_{j\in S_i\subseteq J}T\langle P_j^{m,n}\rangle$
\item For $i\in[2\ell]$, let $g_i$ be the function computing by the $u_i$th gate of $\operatorname{DNF}_{T\langle S_i^m\rangle}$
\item Let $T(x)=D(x,g_1(x),\ldots,g_{2\ell}(x))$
\end{enumerate}
\paragraph{Step4: Describing the $g_i$}
Since $g_1,\ldots,g_{2\ell}$ are components of $g$ and $g=\mathop{\bullet}_{S\in\mathcal S}\operatorname{Eval}\text{-}\operatorname{DNF}_{T\langle S^m\rangle}$, there exist $u_1,\ldots,u_{2\ell}$ and $S_1,\ldots,S_{2\ell}$ such that each $g_i$ is the $u_i$th gate of $\operatorname{DNF}_{T\langle S_i^m\rangle}$ for $i\in[2\ell]$. Moreover, each $u_i\leq2m\log m$ by the trivial upper bound on the number of gates in a canonical DNF.
\paragraph{Step3: Describing the $T\langle S_i^m\rangle$}
Next, we focus on encoding $T\langle S_i^m\rangle$ for some $i$. Since $S_i^m=\bigcup_{j\in S_i}P_j^{m,n}$ (by construction of $S_i^m$), we have (by construction of $T\langle S_i^m\rangle$) that $T\langle S_i^m\rangle=\bigvee_{j\in S_i}T\langle P_j^{m,n}\rangle$. Thus, to compute $T\langle S_i^m\rangle$ for all $i\in[2\ell]$, it suffices to know $S_1,\ldots,S_{2\ell}$ as well as $T\langle P_j^{m,n}\rangle$ for all $j\in J\stackrel{\mathrm{def}}=\bigcup_{i\in[2\ell]}S_i$.
\paragraph{Step2: Describing $T\langle P_j^{m,n}\rangle$ for $j\in J$}
The key to our encoding is to realize that $|J|$ cannot be too large, and thus, we do not need to know $T\langle P_j^{m,n}\rangle$ for all $j\in[n]$. Since $\operatorname{cover}([n],\mathcal S)\geq4\ell$, and $J$ is the union of $2\ell$ sets from $\mathcal S$, it follows that $|J|\leq n-2\ell$. Moreover, we have that
\[n/r\leq\operatorname{cover}([n],\mathcal S)<4\ell+4\]
where the first inequality comes from the sets in $\mathcal S$ having cardinality at most $r$ and the second inequality comes from the definition of $\ell$. Thus,
\[|J|\leq n-2\ell<n-\frac n{2r}+2=n(1-\frac1{2r})+2.\]
Moreover, we can encode $T\langle P_j^{m,n}\rangle$ for $j\in J$ very efficiently. To describe $T\langle P_j^{m,n}\rangle$, it suffices to describe $j$ and then give the list of values of $T\langle P_j^{m,n}\rangle$ on the set $P_j^{m,n}$ in lexicographic order. Since $|P_j^{m,n}|\leq m/n+1$ (essentially by construction), we can encode this list of values by a string $V_j\in\{0,1\}^{m/n+1}$ (where we pad this string with extra zeroes if $|P_j^{m,n}|<m/n+1$).
\paragraph{Step1: Counting the bits in the description}
Now, we count the number of bits in our description of $T$. Describing $J$ requires $n$-bits. For $j\in J$, each partial truth table encoding $(j,V_j)$ requires at most $2\log n+m/n+1$ bits. Using that $|J|\leq n(1-\frac1{2r})$, we get that all these encodings require at most $n(1-\frac1{2r})(2\log n+m/n+1)$ bits. Encoding the $r$-bounded subsets $S_1,\ldots,S_{2\ell}\subseteq[n]$ requires at most $2n\ell\leq n^2$ bits (here we use the fact that $4\ell\leq n$ since $\operatorname{cover}([n],\mathcal S)=4\ell$). Encoding the gate numbers $u_1,\ldots,u_{2\ell}$ requires at most $4\ell\log(2m\log m)=O(n\log n)$ (using that $4\ell\leq n$ and that $m=n^{O(1)}$). Finally, describing a circuit with $\ell$ gates and $(n+2\ell)$ input bits where $4\ell\leq n$ requires $O(n\log n)$ bits. Putting this all together and using that $m=O(n^3)$, we get that $T$ can be described using
\begin{align*}&n+(n(1-\tfrac1{2r})+2)(2\log n+m/n+1)+n^2+O(n\log n)\\&\qquad=(1-\tfrac1{2r})m+O(n^2)=(1-\Omega(1))m\end{align*}
bits. Thus, the number of bad $T$ is upper bounded by $2^{(1-\Omega(1))m}$, so the probability that $T$ is bad is at most $2^{-\Omega(m)}$.
\end{proof}
